% GENERATED FILE. Edit canonical lessons, then run npm run generate:books.
% canonical-source-hash: fnv1a64:d3944239bc730829
% canonical-lessons: IT-C237-avvisare, IT-C237-bruciare, IT-C237-buttare, IT-C237-calcolare, IT-C237-colpire

\chapter{To warn, To burn, To throw away, To calculate, To hit}
\label{ch:it-to-warn-to-burn-to-throw-away-to-calculate-to-hit}
\hlchaptermodality{237}

\begin{chapteropening}
\textbf{By the end of this chapter:} \emph{I can say to warn, to burn, to throw away, to calculate and to hit.}

Complete the last lesson of chapter 237: I can say to warn, to burn, to throw away, to calculate and to hit.
\end{chapteropening}

\section[avvisare]{avvisare — to warn, to let know}

\label{lesson:IT-C237-avvisare}



\begin{quote}
Before the new one: say the Italian for to be bored, then the Italian for to have fun.
\end{quote}

\subsection*{You'll want to know: avvisare}
\textbf{avvisare} — \textquotedblleft{}to warn, to let know\textquotedblright{}.

\begin{cousinweb}
\textbf{avvisare} is the verb of \emph{avviso}, \textquotedblleft{}a notice\textquotedblright{}, which $\leftarrow$ Old French \emph{avis}. (EtymDB 2.1)
\end{cousinweb}

\begin{grammarlens}[title={What you've built}]
One more word for this chapter.
\end{grammarlens}

\subsection*{Your turn}
\begin{itemize}
\raggedright
  \item \emph{Say it:} \emph{avvisare}
  \item \emph{Say it:} \emph{avvisare}, once more
  \item \emph{Say it:} say \emph{avvisare}
  \item \emph{Recall:} say the Italian for to be bored, then the Italian for to have fun, then say \emph{avvisare} again
\end{itemize}

\begin{tcolorbox}[breakable,colback=teal!4,colframe=teal!35!black,title={Before you move on}]
What does avvisare mean? (\textquotedblleft{}To warn, to let know\textquotedblright{}.) Say it once more.
\end{tcolorbox}
\section[bruciare]{bruciare — to burn}

\label{lesson:IT-C237-bruciare}



\begin{quote}
Before the new one: say the Italian for to have fun, then the Italian for to warn.
\end{quote}

\subsection*{You'll want to know: bruciare}
\textbf{bruciare} — \textquotedblleft{}to burn\textquotedblright{}.

\begin{cousinweb}
\textbf{bruciare} — its origin is uncertain. EtymDB 2.1 records the word but not where it came from.
\end{cousinweb}

\begin{grammarlens}[title={What you've built}]
One more word for this chapter.
\end{grammarlens}

\subsection*{Your turn}
\begin{itemize}
\raggedright
  \item \emph{Say it:} \emph{bruciare}
  \item \emph{Say it:} \emph{bruciare}, once more
  \item \emph{Say it:} say \emph{bruciare}
  \item \emph{Recall:} say the Italian for to have fun, then the Italian for to warn, then say \emph{bruciare} again
\end{itemize}

\begin{tcolorbox}[breakable,colback=teal!4,colframe=teal!35!black,title={Before you move on}]
What does bruciare mean? (\textquotedblleft{}To burn\textquotedblright{}.) Say it once more.
\end{tcolorbox}
\section[buttare]{buttare — to throw away}

\label{lesson:IT-C237-buttare}



\begin{quote}
Before the new one: say the Italian for to warn, then the Italian for to burn.
\end{quote}

\subsection*{You'll want to know: buttare}
\textbf{buttare} — \textquotedblleft{}to throw away\textquotedblright{}.

\begin{cousinweb}
\textbf{buttare} $\leftarrow$ Old French \emph{bouter}, \textquotedblleft{}to strike, to push\textquotedblright{} — borrowed. (EtymDB 2.1)
\end{cousinweb}

\begin{grammarlens}[title={What you've built}]
One more word for this chapter.
\end{grammarlens}

\subsection*{Your turn}
\begin{itemize}
\raggedright
  \item \emph{Say it:} \emph{buttare}
  \item \emph{Say it:} \emph{buttare}, once more
  \item \emph{Say it:} say \emph{buttare}
  \item \emph{Recall:} say the Italian for to warn, then the Italian for to burn, then say \emph{buttare} again
\end{itemize}

\begin{tcolorbox}[breakable,colback=teal!4,colframe=teal!35!black,title={Before you move on}]
What does buttare mean? (\textquotedblleft{}To throw away\textquotedblright{}.) Say it once more.
\end{tcolorbox}
\section[calcolare]{calcolare — to calculate}

\label{lesson:IT-C237-calcolare}



\begin{quote}
Before the new one: say the Italian for to burn, then the Italian for to throw away.
\end{quote}

\subsection*{You'll want to know: calcolare}
\textbf{calcolare} — \textquotedblleft{}to calculate\textquotedblright{}.

\begin{cousinweb}
\textbf{calcolare} $\leftarrow$ Latin \emph{calculare}, \textquotedblleft{}to reckon\textquotedblright{} — a learned borrowing. (EtymDB 2.1)
\end{cousinweb}

\begin{grammarlens}[title={What you've built}]
One more word for this chapter.
\end{grammarlens}

\subsection*{Your turn}
\begin{itemize}
\raggedright
  \item \emph{Say it:} \emph{calcolare}
  \item \emph{Say it:} \emph{calcolare}, once more
  \item \emph{Say it:} say \emph{calcolare}
  \item \emph{Recall:} say the Italian for to burn, then the Italian for to throw away, then say \emph{calcolare} again
\end{itemize}

\begin{tcolorbox}[breakable,colback=teal!4,colframe=teal!35!black,title={Before you move on}]
What does calcolare mean? (\textquotedblleft{}To calculate\textquotedblright{}.) Say it once more.
\end{tcolorbox}
\section[colpire]{colpire — to hit}

\label{lesson:IT-C237-colpire}



\begin{quote}
Before the new one: say the Italian for to throw away, then the Italian for to calculate.
\end{quote}

\subsection*{You'll want to know: colpire}
\textbf{colpire} — \textquotedblleft{}to hit\textquotedblright{}.

\begin{cousinweb}
\textbf{colpire} — built from \emph{colpo}, \textquotedblleft{}a blow\textquotedblright{}, with \emph{-ire}. \emph{colpo} $\leftarrow$ Vulgar Latin \emph{colpus}, \textquotedblleft{}a stroke\textquotedblright{}. (EtymDB 2.1)
\end{cousinweb}

\begin{grammarlens}[title={What you've built}]
One more word for this chapter.
\end{grammarlens}

\subsection*{Your turn}
\begin{itemize}
\raggedright
  \item \emph{Say it:} \emph{colpire}
  \item \emph{Say it:} \emph{colpire}, once more
  \item \emph{Say it:} say \emph{colpire}
  \item \emph{Recall:} say the Italian for to throw away, then the Italian for to calculate, then say \emph{colpire} again
\end{itemize}

\begin{tcolorbox}[breakable,colback=teal!4,colframe=teal!35!black,title={Before you move on}]
What does colpire mean? (\textquotedblleft{}To hit\textquotedblright{}.) Say it once more.
\end{tcolorbox}
